% Chapter 1
\chapter{Introducción} % Main chapter title
\label{sec:Introduccion} % For referencing the chapter elsewhere, use

La era digital ha transformado radicalmente la forma en que las personas expresan y comparten sus opiniones. Cada día, millones de usuarios publican reseñas de productos, comentarios en redes sociales y opiniones sobre servicios en plataformas digitales. Este fenómeno ha generado un volumen sin precedentes de texto no estructurado que contiene información valiosa sobre la percepción pública. Las organizaciones modernas dependen cada vez más de esta información para tomar decisiones estratégicas: desde ajustar características de productos hasta gestionar crisis de reputación en tiempo real \cite{liu2015sentiment}.

El análisis de sentimientos se ha consolidado como la técnica fundamental para interpretar esta avalancha de opiniones. Su capacidad para clasificar automáticamente texto según polaridad emocional (positiva, negativa, neutra) permite a empresas y gobiernos monitorear la satisfacción del cliente, detectar tendencias emergentes y responder proactivamente a problemas antes de que escalen. Sin embargo, la efectividad de estos sistemas enfrenta un obstáculo metodológico persistente: el desbalance de clases.

En la mayoría de datasets de sentimientos ampliamente utilizados, la distribución de opiniones dista de ser uniforme. Conjuntos de referencia como Stanford Sentiment Treebank 2, TASS-SEPLN y Coursera Reviews exhiben consistentemente una clase mayoritaria que representa más del 60\% de las muestras \cite{gopali2024applicability}. Este desequilibrio no es meramente un inconveniente técnico. Provoca que los modelos converjan hacia estrategias que favorecen sistemáticamente la clase dominante, sacrificando su capacidad para identificar opiniones minoritarias. El resultado es una degradación pronunciada en métricas como el macro F1-score, comprometiendo la utilidad práctica de estos sistemas en escenarios donde detectar cada tipo de sentimiento es igualmente crítico \cite{gopali2024applicability,whitehouse2023llm}.

\subsection{Enfoques Existentes para el Desbalance de Clases}

La comunidad científica ha desarrollado dos paradigmas principales para abordar este desafío. Las técnicas clásicas de oversampling, ejemplificadas por SMOTE (Synthetic Minority Oversampling Technique) y ADASYN (Adaptive Synthetic Sampling), operan mediante interpolación lineal en espacios vectoriales de embeddings \cite{chawla2002smote,he2008adasyn}. Estos métodos identifican ejemplos de la clase minoritaria y generan muestras sintéticas interpolando entre vecinos cercanos en el espacio de representación. Cuando se combinan con clasificadores basados en modelos transformers como BERT, estas técnicas logran mejoras significativas en F1-score \cite{cloutier2023fine,devlin2018bert}.

Por otro lado, los avances recientes en Large Language Models (LLMs) han abierto una vía alternativa. Modelos como GPT-3 y GPT-4 poseen capacidades generativas que permiten crear texto nuevo directamente a partir de descripciones o ejemplos de la clase minoritaria. La investigación de Suhaeni y Yong demostró que agregar reseñas sintéticas generadas por GPT-3 al corpus Coursera Reviews produjo incrementos de 12.7\% en exactitud media y hasta 13 puntos porcentuales en macro F1-score \cite{suhaeni2023mitigating}. Estos modelos generan texto gramaticalmente correcto que preserva naturalidad lingüística, superando la limitación fundamental de los métodos de interpolación: la falta de correspondencia entre vectores interpolados y texto lingüísticamente válido \cite{whitehouse2023llm}.

Cada paradigma presenta fortalezas complementarias. SMOTE y ADASYN proporcionan guía geométrica explícita sobre dónde posicionar muestras sintéticas en el espacio de embeddings, identificando estratégicamente regiones subrepresentadas. Los LLMs, en contraste, generan texto coherente y diverso, pero operan sin información sobre la estructura geométrica del espacio de representación. Esta complementariedad ha motivado el surgimiento de metodologías híbridas emergentes.

\subsection{El Paradigma Híbrido Emergente}

Arquitecturas recientes han comenzado a integrar técnicas clásicas con capacidades generativas de LLMs. SMOTExT propone utilizar SMOTE para identificar puntos de interpolación en el espacio de embeddings, luego emplear un LLM para generar texto correspondiente a esas ubicaciones geométricas \cite{smotext2025arxiv}. De manera similar, ImbLLM explora la capacidad de LLMs para generar muestras sintéticas guiadas por información estructural del dataset \cite{imbllm2024arxiv}. Estos trabajos han establecido las bases conceptuales del paradigma híbrido, demostrando que entrenar exclusivamente con datos sintéticos puede lograr rendimientos cercanos a baselines con datos reales, con degradaciones de aproximadamente 3.2\% en algunos casos.

A pesar de estos avances prometedores, el análisis de la literatura revela que los enfoques híbridos actuales aún presentan limitaciones metodológicas que restringen su efectividad. Estas limitaciones se detallan en la sección de Definición del Problema y constituyen la motivación principal de esta investigación.

\subsection{Contribución de Esta Investigación}

Esta tesis propone una arquitectura híbrida avanzada que integra cuatro innovaciones metodológicas diseñadas para superar las limitaciones identificadas en enfoques previos. Primero, implementa una selección de soporte ponderada que utiliza un número mayor de vecinos contextualmente relevantes para proporcionar información más rica al LLM durante la generación. Segundo, incorpora un sistema de filtrado multi-nivel adaptativo que valida la calidad de las muestras sintéticas mediante criterios complementarios aplicados en cascada. Tercero, aplica ponderación basada en confianza geométrica, donde cada muestra sintética recibe peso proporcional a indicadores de su contexto local en el espacio de embeddings. Finalmente, el diseño enfatiza la eficiencia computacional mediante generación directa con LLMs pre-entrenados sin requerir fine-tuning.

El objetivo principal es evaluar sistemáticamente, mediante estudios de ablación rigurosos en múltiples datasets con diversos grados de desbalance, las contribuciones individuales de cada componente y la efectividad del sistema integrado en tareas de clasificación de sentimientos. Esta investigación busca proporcionar un protocolo reproducible que guíe la selección de técnicas según características del dataset y restricciones operacionales, contribuyendo tanto al avance científico en el equilibrio de clases como a la aplicabilidad práctica en sistemas de producción que requieren predicciones confiables en todas las categorías de sentimiento.